% Estilo comum das figuras de listas (incluído pelas figuras dos capítulos 3 e 4).
\usepackage{tikz}
\usetikzlibrary{arrows.meta, positioning, shapes.multipart, calc}
\definecolor{navy}{RGB}{23,45,80}
\definecolor{laranja}{RGB}{250,110,29}
\definecolor{cinzalogo}{RGB}{107,114,126}
\tikzset{
  no/.style={rectangle split, rectangle split horizontal, rectangle split parts=2,
             draw=navy, line width=0.9pt, minimum height=7mm,
             rectangle split part fill={laranja!15, white}, font=\sffamily\small},
  no3/.style={rectangle split, rectangle split horizontal, rectangle split parts=3,
              draw=navy, line width=0.9pt, minimum height=7mm,
              rectangle split part fill={white, laranja!15, white}, font=\sffamily\small},
  seta/.style={-{Stealth[length=2.2mm]}, navy, line width=0.9pt},
  setanova/.style={-{Stealth[length=2.2mm]}, laranja, line width=1.3pt},
  rotulo/.style={font=\sffamily\small\bfseries, text=navy},
  nulo/.style={font=\sffamily\scriptsize, text=cinzalogo},
  passo/.style={font=\sffamily\small, text=cinzalogo, anchor=west},
  legenda/.style={font=\sffamily\scriptsize, text=cinzalogo},
}
% seta que sai do meio da caixa do ponteiro (última parte do nó), na horizontal
\newcommand{\liga}[3][seta]{\draw[#1] ([xshift=-2.3mm]#2.east) -- #3;}
